\documentclass[12pt]{article}

%importamos los paquetes
\usepackage[utf8]{inputenc} %Configura caracteres especiales
\usepackage[spanish]{babel} %Configura en español los valores predet.
\usepackage{amsmath} %Para el manejo de formulas
\usepackage{array} %Mejor manejo de las tablas
\usepackage{geometry} %Margenes
\geometry{margin=2.5cm} %Modifica el margen de la hoja

\newcommand{\minombre}{brandon alexis avalos santillan}
\newcommand{\nombrepractica}{Practica 4: Tablas, entornos y comandos}
\newcommand{\universidad}{UASLP-FEPZM}

\newenvironment{ejercicio}
{
 \vspace{0.3cm}
\noindent\textbf{Ejercicio:}
\begin{itshape} %Modifica a italika el texto del ejercicio
}
{
   \end{itshape}
\vspace{0.3cm}
}

\begin{document}
\begin{center}
    {\LARGE \nombrepractica}\\[0.3cm]
\vspace{10cm}
{\Large  \minombre}\\
\vspace{10cm}
\universidad\\
{\large \today}
\end{center}

\section{Introduccion}
En esta practica aprenderemos a crear tablas y a usar \textbf{environments}, incluyendo environments personalizados

\section{Uso de Tablas}
\begin{table}[h]
    \centering
    \begin{tabular}{c|c|c}
         %\hline
         Alumno&Materia&Calificacion  \\ 
         \hline
         Pedro&Matematicas&8\\
         Maria&Español&10\\
    \end{tabular}
    \caption{Calificaciones de alumnos}
    \label{Calificaciones}
\end{table}
\subsection{Tablas con columnas combinadas}
\begin{table}[h]
    \centering
    \begin{tabular}{c|c|c}
    \hline
    \multicolumn{3}{|c|}{Horario}\\ \hline
    Dia&Hora&Materia\\ \hline
    Lunes&10 a 11&Matematicas\\ \hline
    Martes&12 a 13&Fisica\\ \hline
    
    \end{tabular}
    \caption{Horario}
    \label{Horarario}

\section{Environment Personalizado}
Ahora usaremos nuestro \textbf{environment}

\begin{ejercicio}
    Crear una tabla con almenos 4 filas y 3 columnas, rellenala con contenido afin a un tema en especifico
\end{ejercicio}



\end{table}

\end{document}
